\include{preambula_questions}
\usepackage[most]{tcolorbox}
\setlength{\parindent}{0pt}
\setlength{\parskip}{6pt}

\begin{document}
\title{Генетический алгоритм}
\maketitle
\section{Эволюционные алгоритмы}
Эволюционные алгоритмы — это семейство алгоритмов оптимизации, основанных на идее биологической эволюции и естественного отбора Чарльза Дарвина. Генетические алгоритмы являются самым известным и широко используемым подклассом этого семейства. Моделирование эволюции с использованием идей эволюционных алгоритмов и искусственной жизни началось с работы Нильса Aall Barricelli в 1960-х, и было продлено Алексом Фрейзером, который опубликовал ряд работ по моделированию искусственного отбора. Эволюционные алгоритмы стали общепризнанным методом оптимизации в результате работ Инго Rechenberg в 1960-х и начале 1970-х, который использовал их для решения сложных инженерных задач. \\
Идейно суть метода заключается в упрощенном моделировании процессов эволюции для поиска оптимального (или достаточно хорошего) решения сложной задачи. Вместо того чтобы писать алгоритм для поиска точного ответа, который гарантированно сойдётся (но не факт, что мы доживем до этого момента), мы создаем «популяцию» возможных решений и заставляем их эволюционировать, пока не получим приемлемый результат. \\
Помимо классического генетического алгоритма, в класс эволюционных алгоритмов входят: \\
1. Генетическое программирование: эволюционируют не параметры, а сами компьютерные программы (в виде деревьев выражений). \\
2. Эволюционные стратегии: акцент делается на мутациях и адаптации параметров мутации прямо во время работы (например, CMA-ES). \\
3. Метаэвристики на основе роевого интеллекта: алгоритм роя частиц (PSO), муравьиные алгоритмы (Ant Colony) — вдохновлены не генами, а социальным поведением стайных животных. \\

Поговорим о мотивировке использования алгоритмов этого класса. \\
Глубокое обучение требует огромных, хорошо размеченных датасетов и мощных GPU, но архитектуры и методы обучения стандартизированы. Эволюционные алгоритмы — это скорее «инструмент плачущего инженера», который применяется, когда классические методы бессильны. \\
Причины использовать ГА: \\
1. Сложность задач: задачи часто относятся к классу NP-трудных (например, задача коммивояжера с тысячами городов). Идеального математического решения за разумное время не существует. \\
2. Нестабильность требований: задача постоянно меняется. Добавились новые ограничения (например, мы хотим написать оптимизатор маршрута курьера в городе, а бизнес добавил требования: девушки не могут унести больше 10кг за раз, а потом оказалось, что каждый второй парень не работает по вторникам), и алгоритм, который точно сойдётся, нужно переписывать с нуля. ГА адаптируется просто добавлением штрафа в функцию приспособленности. \\
3. Жесткие временные рамки: нет времени на разработку идеального решения. Нужно сделать «хоть как-то, но чтобы работало». ГА позволяет получить приближенное, но рабочее решение без гарантий абсолютного оптимума. \\

Некоторые примеры использования ГА: \\
1. Поиск новых материалов (Программа USPEX): алгоритм предсказывает структуру кристаллов при гигантских давлениях (таких, как в ядре Земли). Так были открыты соединения $Na_3Cl$ и $NaCl_7$, которые нарушают классическую школьную химию, но существуют физически. \\
2. Генерация тестов: ГА может эволюционировать не решение, а саму проблему. Алгоритм генерирует такие входные данные, на которых программа разработчика будет работать максимально долго или потреблять максимум памяти.
3. Генетическое улучшение кода (Genetic Improvement): использование ГА прямо на исходном коде программ. Алгоритм случайным образом меняет, удаляет или копирует строки кода, а затем прогоняет тесты. Так удалось ускорить программу секвенирования генома человека в 77 раз просто за счет неочевидных машинных оптимизаций.
\newpage
\section{Основные понятия и базовый цикл}
Особь (индивид) — одно конкретное решение задачи. \\
Популяция — множество особей, существующих в данный момент времени. \\
Приспособленность (целевая функция, fitness) — числовая оценка того, насколько хорошо данная особь решает задачу. \\
Ген/Хромосома — параметры, из которых состоит решение (например, город в маршруте или значение конкретной переменной). \\
Любой генетический алгоритм работает по определенному циклу, который повторяется из поколения в поколение: \\
1. \textbf{Генерация начальной популяции.} Создается набор случайных (или псевдослучайных) решений-кандидатов. В ГА такими объектами выступают вектора («генотипы») генов, где каждый ген может быть битом, числом или неким другим объектом. \\
2. \textbf{Оценка приспособленности.} Каждое решение прогоняется через целевую функцию. \\
3. \textbf{Селекция.} Из популяции выбираются особи для создания потомства. Чем лучше решение, тем больше у него шансов быть выбранным. \\
4. \textbf{Скрещивание.} Две (или более) родительские особи обмениваются своими генами, порождая новые решения, которые комбинируют признаки родителей. Главное требование к размножению — чтобы потомок или потомки имели возможность унаследовать черты обоих родителей, «смешав» их каким-либо способом.\\
5. \textbf{Мутация.} С небольшой вероятностью в новые решения вносятся случайные изменения. Мутация не дает алгоритму застрять в локальном минимуме. В ГА оператор мутации может быть реализован простым добавлением нормально распределенной случайной величины к каждой компоненте вектора. \\
6. \textbf{Отбор для выживания.} Из старого поколения и новых потомков формируется новая популяция. Худшие решения «умирают», лучшие переходят в следующее поколение (или становятся его основой). \\
7. \textbf{Остановка.} Цикл повторяется до выполнения критерия останова.

\section{Дискретная оптимизация}
Любая задача оптимизации формулируется как поиск такого вектора параметров $x^*$ в заданном пространстве поиска $\mathcal{X}$, при котором целевая функция $f(x)$ достигает своего экстремума (для определенности — максимума): \\
$x^* = \arg\max_{x \in \mathcal{X}} f(x)$

Ключевой проблемой является наличие локальных оптимумов. Точка $x$ является строгим локальным максимумом, если для всех точек $x'$ в некоторой её окрестности выполняется $f(x) > f(x')$. Локальные алгоритмы по своей природе сходятся к локальному оптимуму и не способны его покинуть, так как любой следующий шаг в окрестности приведет к ухудшению значения целевой функции.

В непрерывных пространствах ($\mathcal{X} \subset \mathbb{R}^n$) для определения направления поиска используются производные целевой функции. \\

1. Градиентный спуск. \\
Базовый метод оптимизации первого порядка.
Алгоритм итеративно обновляет текущее решение, сдвигаясь в направлении вектора градиента. Правило обновления: $x_{k+1} = x_k + \alpha \nabla f(x_k)$, где $\alpha$ — размер шага. \\
Алгоритм останавливается, когда достигает стационарной точки . Эта точка может быть локальным максимумом, минимумом или седловой точкой. Алгоритм принципиально не способен преодолевать зоны с убыванием целевой функции для поиска более оптимальных областей. Кроме того, метод требует дифференцируемости целевой функции. \\
2. Метод Ньютона. \\
Метод оптимизации второго порядка, использующий информацию о кривизне поверхности целевой функции. \\
Метод использует Гессиан, $H$. Шаг вычисляется с учетом кривизны: $x_{k+1} = x_k + \alpha H^{-1} \nabla f(x_k)$. Это позволяет алгоритму совершать более длинные шаги на пологих участках и более короткие — на участках с высокой кривизной.
Метод обладает квадратичной скоростью сходимости на гладких унимодальных функциях. Однако вычислительная сложность расчета и обращения матрицы Гессе на каждой итерации составляет $O(n^3)$, что делает метод неприменимым для задач высокой размерности. Как и градиентный спуск, метод Ньютона является строгим локальным оптимизатором и не способен выходить из локальных экстремумов.

При переходе к дискретной оптимизации (комбинаторной оптимизации) решения представляют собой структуры данных фиксированной или переменной длины (графы, перестановки, битовые строки). В таких пространствах $\mathcal{X}$ понятие производной не определено. Математический аппарат непрерывной оптимизации заменяется концепцией локального изменения (мутации): \\
Вместо градиента вводится оператор мутации $m: \mathcal{X} \to \mathcal{X}$, представляющий собой минимально возможное изменение структуры решения-кандидата. \\
Применение всех возможных мутаций к текущей точке $x$ формирует её окрестность $\mathcal{N}(x)$. \\
Процесс оптимизации формализуется как направленный обход графа, где вершинами являются решения-кандидаты, а ребрами — возможные локальные изменения. Методы, реализующие такой обход, классифицируются как алгоритмы локального поиска, или Hill Climbers. \\
Посмотрим на некоторые алгоритмы локального поиска, для этого рассмотрим следующую задачу: имеем алфавит конечной длины и строку длины N. Требуется максимизировать сумму ASCII кодов символов входящих в строку. \\
Пространство поиска $\mathcal{X}$: множество строк фиксированной длины $N$, элементы которых принадлежат алфавиту $\Sigma$. В данном случае $\Sigma$ — множество заглавных букв латинского алфавита, $|\Sigma| = 26$. \\
Целевая функция $f(x)$: аддитивная функция, вычисляющая сумму ASCII-кодов всех символов вектора $x$. \\
$$f(x) = \sum_{i=1}^N \text{ASCII}(x_i)$$
Глобальный оптимум $x^*$: очевиден аналитически. Строка, состоящая исключительно из символов с максимальным ASCII-кодом в заданном алфавите (строка из $N$ символов «Z»). Для удобства вычислений сопоставим символам их числовые значения от $1$ (буква 'A') до $S = |\Sigma|$ (буква 'Z'). \\
Начальное состояние $x_0$: для наглядности анализа стартовой точкой принимается строка с минимальным значением целевой функции (строка из $N$ символов «A»). \\
Оператор локальной мутации $M$: функция, заменяющая ровно один символ $x_i$ на любой другой символ из алфавита $\Sigma$. Размер окрестности при таком операторе составляет $|\mathcal{N}(x)| = N \cdot (|\Sigma| - 1)$.
Попробуем оценить математическое ожидание количества вычислений целевой функции, необходимых для достижения глобального оптимума $x^*$ из точки $x_0$, используя различные стратегии: \\
\textbf{1. First Ascent Hill Climber} \\
Алгоритм осуществляет детерминированный последовательный перебор списка возможных мутаций $M$. Итерация завершается, как только найдена мутация $m(x) \to y$, удовлетворяющая строгому условию $f(y) > f(x)$. Найденное решение $y$ становится текущим, и перебор списка мутаций начинается заново. \\
Пусть список мутаций $M$ упорядочен (например, сначала все замены для 1-го символа, затем для 2-го и т.д.). Размер списка $|M| = N \cdot S$. Чтобы полностью оптимизировать строку, нам нужно суммарно сделать $N \times (S-1)$ успешных единичных улучшений (так как алгоритм наткнется на ближайшую улучшающую букву, инкрементируя значение). С каждым новым успехом количество уже оптимизированных позиций, которые алгоритм должен проверить и отвергнуть при старте с начала списка, растет. На $k$-м успешном шаге алгоритму придется в среднем просмотреть $\mathcal{O}(k)$ отвергаемых мутаций, прежде чем он дойдет до той, которая даст улучшение. Общее количество вычислений функции представляет собой сумму арифметической прогрессии от 1 до максимального числа проверок (которое пропорционально длине списка $N \cdot S$):
    $$E[T_{FAHC}] \approx \sum_{k=1}^{N \cdot S} k = \frac{(N \cdot S)(N \cdot S + 1)}{2} \approx \frac{S^2 N^2}{2}$$
\textbf{2. Steepest Ascent Hill Climber} \\
На каждой итерации алгоритм вычисляет целевую функцию для каждого элемента окрестности $\mathcal{N}(x)$. Осуществляется переход в точку $y$, обеспечивающую наибольшее приращение функции. Для достижения строки «Z...Z» из «A...A» требуется совершить ровно $N$ успешных макрошагов. На каждом таком шаге необходимо полностью вычислить окрестность, размер которой пропорционален $N \cdot |\Sigma|$.
Итоговая сложность: $O(|\Sigma| \cdot N^2)$. \\

\textbf{3. Randomized Local Search} \\
Алгоритм выбирает случайную мутацию из $M$ (выбирает случайный индекс от $1$ до $N$ и случайный символ из $\Sigma$). Мутация принимается при выполнении нестрогого условия: $f(y) \ge f(x)$. 
На начальных этапах (когда строка состоит из «A») почти любая случайная мутация ведет к улучшению. Однако по мере приближения к вектору $x^*$ вероятность успешной генерации мутации $\left(p \approx \frac{1}{N \cdot |\Sigma|}\right)$ падает. \\
Пусть текущее значение целевой функции равно $f(x)$. Максимально возможное значение равно $F_{max} = N \cdot S$. Назовем «дистанцией до оптимума» величину $d = F_{max} - f(x)$. В начале $d = N \cdot S - N \approx N \cdot S$. \\
Посмотрим на текущую строку $x$. Для любого символа $x_i$ существует ровно $(S - x_i)$ мутаций, которые заменят его на символ с большим кодом. Просуммируем количество всех строго улучшающих мутаций для всей строки: \\
   $$\ \sum_{i=1}^N (S - x_i) = N \cdot S - \sum_{i=1}^N x_i = N \cdot S - f(x) = d$$
То есть количество строго улучшающих мутаций в точности равно дистанции до оптимума $d$. На каждом шаге RMHC выбирает одну мутацию из $N \cdot S$ возможных абсолютно случайно. Следовательно, вероятность выбрать строго улучшающую мутацию при текущей дистанции $d$ равна:
    $$P_d = \frac{d}{N \cdot S}$$
По свойствам геометрического распределения, ожидаемое количество вычислений функции до того, как мы выберем улучшающую мутацию и дистанция $d$ сократится хотя бы на 1, составляет:
    $$E[T_d] = \frac{1}{P_d} = \frac{N \cdot S}{d}$$
Чтобы дойти до оптимума, мы должны пройти все дистанции от начальной $d \approx N \cdot S$ до $d=1$. Суммируем математические ожидания времени пребывания на каждом этапе:
    $$E[T_{total}] \le \sum_{d=1}^{N \cdot S} \frac{N \cdot S}{d} = N \cdot S \sum_{d=1}^{N \cdot S} \frac{1}{d} \approx N \cdot S \cdot \ln(N \cdot S)$$
Однако есть существенный недостаток: вероятность покинуть строгий локальный оптимум равна нулю. Отсюда естественным образом возникает вопрос: как модифицировать локальный алгоритм, чтобы превратить его в глобальный оптимизатор?

\section{Глобальная оптимизация}
Глобальным оптимизатором называется алгоритм, который имеет ненулевую вероятность достичь глобального оптимума из любой начальной точки пространства за конечное число шагов. \\

\textbf{1. Эволюционный алгоритм $(1+1)$-EA: Стандартная мутация} \\
Простейшим шагом от локального поиска RMHC к эволюционным алгоритмам является алгоритм $(1+1)$-EA. Запись $(1+1)$ означает, что один родитель генерирует одного потомка, после чего происходит отбор. Чтобы алгоритм стал глобальным, необходимо изменить оператор мутации. Вместо замены ровно одного случайного символа применяется следующий оператор мутации: \\
Алгоритм проходит по каждому из $N$ символов строки. Каждый символ подвергается мутации (замене) независимо от остальных с фиксированной вероятностью $p = 1/N$. \\
Ожидаемое количество измененных символов (согласно свойствам биномиального распределения $\mathcal{B}(N, 1/N)$) равно $N \cdot (1/N) = 1$. То есть в среднем алгоритм ведет себя точно так же, как эффективный локальный поиск RMHC, не теряя скорости на простых унимодальных участках. \\
Принципиальное отличие в том, что теперь существует ненулевая вероятность изменить 2, 3 или вообще все $N$ символов за одну итерацию. Следовательно, алгоритм теоретически способен перепрыгнуть любой локальный оптимум. \\
Однако вероятность изменить $k$ символов одновременно убывает экспоненциально при росте $k$ - в этом существенный недостаток.

\textbf{2. Fast $(1+1)$-EA: Мутация по степенному закону} \\
Чтобы решить проблему экспоненциального падения вероятности больших прыжков предлагается следующая модификация (метод предложен совсем недавно: в 2014 году). Суть модификации заключается в отказе от биномиального распределения в пользу степенного закона. \\
На каждой итерации алгоритм сначала выбирает количество символов $c \in [1..N]$, которые нужно изменить. Это число $c$ выбирается с вероятностью, пропорциональной $c^{-\beta}$ (где $\beta$ — константа, обычно $\beta \in (1, 2)$). Выбрав $c$, алгоритм случайно меняет ровно $c$ символов. \\
Степенное распределение обладает так называемым «тяжелым хвостом». Это означает, что вероятность выбрать $c=1$ всё еще самая высокая (алгоритм в среднем остается локальным поиском). Однако вероятность выбрать большое значение (например, $c=N/2$) убывает полиномиально, а не экспоненциально. 

Алгоритм Fast $(1+1)$ тратит на поиск оптимума в простых задачах почти столько же времени (асимптотически $O(N \log N)$ с чуть большей константой). Но на сложных задачах он показывает на порядки лучшую производительность, так как способен за адекватное время перепрыгивать широкие долины между локальными оптимумами. \\

\textbf{3. Увеличение популяции: $(1+\lambda)$ и $(1,\lambda)$ Эволюционные Стратегии} \\
До этого момента рассматривался базовый алгоритм $(1+1)$-EA, где единственная текущая точкапорождала ровно одного кандидата на каждой итерации. Теперь пусть от каждого родителя будет $\lambda$ потомков. В зависимости от механизма селекции, возникают две принципиально разные стратегии. \\
1. Стратегия с плюсом: $(1 + \lambda)$-EA (Элитизм) \\
На $k$-й итерации родитель $x_k$ порождает $\lambda$ независимых потомков $y_1, y_2, \dots, y_\lambda$ с помощью оператора мутации. На этап отбора выставляется пул из $\lambda + 1$ особей: все потомки и сам родитель. Следующим родителем $x_{k+1}$ становится особь с максимальным значением фитнеса из этого пула. Если лучший потомок хуже родителя, родитель переходит на следующий шаг неизменным. Значение целевой функции $f(x)$ в последовательности поколений монотонно не убывает: $f(x_{k+1}) \ge f(x_k)$. \\
Из-за строгого элитизма $(1 + \lambda)$-EA, по своей сути, остается локальным оптимизатором. Если родитель $x_k$ оказался в строгом локальном максимуме, все его $\lambda$ потомков с подавляющей вероятностью будут иметь меньший фитнес. В результате родитель будет бесконечно выживать, и алгоритм стагнирует. \\
2. Стратегия с запятой: $(1, \lambda)$-EA (Принудительная деградация) \\
Родитель порождает $\lambda$ потомков. На этап отбора выставляется пул только из $\lambda$ потомков. Родитель безусловно уничтожается. Выбирается лучший из потомков: $x_{k+1} = \arg\max_{y_i} f(y_i)$.
Отсутствие элитизма позволяет алгоритму совершать переходы, при которых $f(x_{k+1}) < f(x_k)$. \\
Предположим, родитель достиг локального пика. Все сгенерированные потомки лежат ниже по склону (имеют худший фитнес). Поскольку родитель уничтожается, алгоритм вынужден выбрать наилучшего из худших потомков. Таким образом, алгоритм целенаправленно спускается с локального холма. Если параметр $\lambda$ достаточно велик (пространство исследовано плотно), алгоритм выберет оптимальную траекторию спуска, что позволит ему пересечь бассейн локального минимума и начать восхождение на новый, потенциально более высокий, локальный оптимум.



\textbf{Перезапуск алгоритмов} \\
Если алгоритмы локального поиска (FAHC, SAHC, RMHC) неизбежно застревают в локальных оптимумах, эту проблему можно решить экстенсивным методом — внедрением глобального цикла перезапусков.
1. Инициализировать глобальную переменную $b$ (best) случайным решением. \\
2. В цикле (пока не исчерпан лимит времени или итераций): \\
3. Сгенерировать новую, случайную начальную точку $x$, равномерно распределенную в пространстве поиска $\mathcal{X}$. \\
4. Запустить из точки $x$ процедуру LocalSearch(x), которая работает до строгой сходимости в локальный оптимум $y$.
Если найденный локальный оптимум лучше глобального рекордсмена ($f(y) > f(b)$), обновить рекордсмен: $b = y$. Возвращаем $b$. \\

Пространство поиска $\mathcal{X}$ можно разбить на «бассейны притяжения» (basins of attraction). Бассейн притяжения оптимума $y$ — это множество всех начальных точек $x$, стартуя из которых детерминированный локальный поиск неизбежно сойдется к $y$. Поскольку глобальный оптимум $y^*$ тоже имеет свой бассейн притяжения размером $V > 0$, вероятность того, что абсолютно случайная точка $x$ попадет в этот бассейн, строго больше нуля. При стремлении числа перезапусков к бесконечности, вероятность нахождения глобального оптимума сходится к 1. \\

\textbf{Снова про отжиг} \\
Главный недостаток отжига - подбор начальной температуры и графика ее понижения для каждой новой задачи. Это отдельное шаманство, требующее ручной настройки гиперпараметров, для чего нужно много довольно знать про целевую функцию. Поэтому на практике исследователи и инженеры часто предпочитают более современные эволюционные алгоритмы (например, CMA-ES для непрерывных или Fast EA для дискретных сред), которые обладают встроенными механизмами самоадаптации и требуют меньшего ручного вмешательства.\\

\textbf{Tabu Search} \\
Вводится структура данных $T$ (Tabu list — список запретов) и параметр $L$ (максимальный размер этого списка). \\
1. Алгоритм генерирует случайную мутацию $y = m(x)$. \\
2. Проверяется два условия для принятия нового решения $y$: \\
$f(y) \ge f(x)$ (стандартное условие локального поиска: решение не должно быть хуже). \\
$y \notin T$ (решение не должно находиться в списке запретов). \\
Если оба условия выполнены, $y$ принимается как текущее решение и добавляется в множество $T$. Так как память ограничена, если размер $|T| > L$, из списка удаляется самый старый элемент.
Однако вариант с запоминанием состояний (самих точек $x$) в дискретных пространствах работает крайне плохо. В задачах комбинаторной оптимизации пространство поиска огромно, вероятность случайно сгенерировать мутацию и попасть в точности в то же самое состояние, в котором алгоритм был 10 шагов назад, околонулевая. Следовательно, проверка $y \notin T$ будет почти всегда выдавать True, и список запретов просто не будет работать. Это приводит к необходимости модификации метода. \\

Чтобы метод заработал в многомерных дискретных пространствах, подход меняется: табуируются не состояния, а сами операции мутации. \\
Вместо проверки наличия состояния $y$ в памяти, алгоритм модифицирует сам генератор мутаций: теперь $T$ хранит не векторы $x$, а примененные операторы $m$. Логика в следующем: главная задача алгоритма в локальном оптимуме — заставить его спуститься вниз по склону. Если мы сделали шаг вниз (вынужденно приняли ухудшающую мутацию, чтобы выйти из ямы), самым выгодным шагом на следующей итерации с точки зрения градиента будет шаг обратно на вершину. Чтобы этого не произошло, мы заносим в список $T$ обратную мутацию. \\
В строгом математическом смысле Tabu Search с конечным списком памяти $L$ не гарантирует нахождения глобального оптимума при $t \to \infty$. Он может попадать в сложные циклы, длина которых превышает размер списка запретов $L$.
Также алгоритм крайне чувствителен к размеру списка $L$. Если $L$ слишком мал, алгоритм зациклится. Если $L$ слишком велик, алгоритм может заблокировать пути к действительно хорошим (еще не исследованным) решениям, так как запретит слишком много полезных мутаций. \\

\textbf{Iterated Local Search} \\
Классические алгоритмы локального поиска (FAHC, SAHC, RMHC) осуществляют навигацию в полном дискретном пространстве решений $\mathcal{X}$, переходя от точки к точке, пока не достигнут локального оптимума. Их главный недостаток — необратимая остановка в этой точке.

Random Restarts решает эту проблему путем генерации новой случайной точки в пространстве $\mathcal{X}$. Однако этот подход обладает амнезией: он игнорирует топологию целевой функции. В задачах комбинаторной оптимизации хорошие решения часто сгруппированы в определенных областях пространства. Абсолютно случайный рестарт выбрасывает алгоритм из этой перспективной зоны, заставляя тратить вычисления на исследование заведомо низкокачественных областей. \\
Мы хотим, чтобы алгоритм осуществлял поиск не в полном пространстве всех возможных решений $\mathcal{X}$, а в редуцированном подпространстве $\mathcal{X}^* \subset \mathcal{X}$, которое состоит исключительно из локальных оптимумов. ILS трансформирует процесс оптимизации в направленное случайное блуждание по локальным экстремумам. \\
Пусть $f: \mathcal{X} \to \mathbb{R}$ — целевая функция, которую необходимо максимизировать. \\
1.  Генерация начального решения: $x_0 \in \mathcal{X}$. \\
2.  Локальная оптимизация: $x^*_0 = \text{LocalSearch}(x_0)$. Точка $x^*_0 \in \mathcal{X}^*$ становится базовой (текущим локальным оптимумом). \\
3.  Итеративный цикл: \\
Возмущение: $x' = \text{LargeMutation}(x^*_k)$. \\
Локальный спуск: $x^{*\prime} = \text{LocalSearch}(x')$. \\
Критерий принятия: $x^*_{k+1} = \text{NewBase}(x^*_k, x^{*\prime})$. \\

1. Оператор LocalSearch \\
Отображение $\text{LS}: \mathcal{X} \to \mathcal{X}^*$. Оператор должен быть максимально вычислительно эффективным. Его единственная задача — "скатить"\, решение на дно текущего бассейна.

2. Оператор LargeMutation \\
Вывод решения из состояния локального оптимума путем значительного изменения его структуры. Эффективность алгоритма зависит от метрики между $x^*_k$ и возмущенной точкой $x'$. Возможны три сценария: \\
1.  Недостаточное возмущение: Если $x'$ слишком близка к $x^*_k$, то $x'$ останется внутри текущего бассейна притяжения $B(x^*_k)$. В этом случае последующий вызов $\text{LocalSearch}(x')$ детерминированно вернет алгоритм обратно в точку $x^*_k$. Произойдет холостая итерация, алгоритм зациклится (вероятность выхода из оптимума равна 0). \\
2.  Избыточное возмущение: Если $x'$ слишком далека от $x^*_k$, возмущение вырождается в генерацию абсолютно случайной точки в $\mathcal{X}$. В этом случае ILS полностью теряет свою эффективность и деградирует до метода слепых случайных перезапусков (Random Restarts), теряя информацию о найденной "глобальной долине". \\
3.  Оптимальное возмущение: Величина изменения $\Delta x$ должна быть минимально необходимой для того, чтобы точка $x'$ пересекла границу текущего бассейна притяжения $B(x^*_k)$ и оказалась в **соседнем** бассейне притяжения $B(x^{*\prime})$. Таким образом алгоритм перепрыгивает из одного локального оптимума в ближайший к нему другой локальный оптимум.

3. Оператор NewBase \\
Определяет, станет ли новый локальный оптимум $x^{*\prime}$ базовой точкой для следующей итерации, или алгоритм отвергнет его и вернется к $x^*_k$. Данный оператор задает стратегию навигации в надпространстве $\mathcal{X}^*$.
Варианты реализации (модульность метаэвристик):
    1.  Строгий (Жадный) выбор* $x^{*\prime}$ принимается только если $f(x^{*\prime}) > f(x^*_k)$. В этом случае ILS работает как алгоритм FAHC, но не на уровне обычных решений, а на макро-уровне локальных оптимумов. Недостаток: алгоритм может застрять в "макро-локальном оптимуме".
    2.  Случайное блуждание: $x^{*\prime}$ принимается всегда. Обеспечивает максимальную глобальность, но страдает от низкой скорости сходимости.
    3.  Вероятностный выбор (Simulated Annealing):* Принятие худшего локального оптимума допускается с вероятностью, зависящей от разницы фитнесов $\Delta f$ и параметра "макро-температуры". Это классический пример, о котором упоминает лектор: встраивание логики имитации отжига внутрь ILS для достижения идеального баланса между разведкой и уточнением.

Мощь Iterated Local Search заключается в радикальном уменьшении кардинальности пространства поиска. Вместо того чтобы оценивать асимптотику на множестве $|\mathcal{X}|$, ILS оперирует на множестве $|\mathcal{X}^*|$, где $|\mathcal{X}^*| \ll |\mathcal{X}|$. За счет сохранения базовой точки и использования контролируемого возмущения, алгоритм исследует коррелированные локальные оптимумы. Если в задаче существует сильная автокорреляция ландшафта, ILS будет целенаправленно дрейфовать в сторону глобального экстремума, избегая экспоненциальных затрат времени, присущих методам слепого перебора или стандартным мутациям с независимыми битовыми вероятностями.


\end{document}